\subsection{变分问题的强制性}
\label{sec:ch05-variational-coercivity}
\setcounter{thmcount}{0}
\setcounter{equation}{0}

为应用第~\ref{chap:ch02-variational-formulation} 章建立的理论, 必须说明前两节引入的变分形式 $a(\cdot,\cdot)$ 在相应空间 $V$ 上具有强制性; 它在整个 $H^1(\Omega)$ 上有界是显然的. 纯 Neumann 边界条件情形的强制性来自第~\ref{chap:ch04-polynomial-approximation} 章的逼近理论. 若 $\Omega$ 是有界区域, 可写成有限个关于某个球为星形的区域之并, 则
\MBSourceItem{1}
\begin{equation}
\MathBookFormulaID{formula-ch05-best-constant-poincare}
\label{eq:ch05-best-constant-poincare}
\inf_{r\in\mathbb R}\|v-r\|_{L^2(\Omega)}
\leq C_\Omega|v|_{H^1(\Omega)}
\quad\forall v\in H^1(\Omega).
\end{equation}
又有
\[
\MathBookFormulaID{formula-ch05-mean-is-best-constant}
\inf_{r\in\mathbb R}\|v-r\|_{L^2(\Omega)}
=\|v-\overline v\|_{L^2(\Omega)},
\]
其中 $\overline v$ 是式~\eqref{eq:ch05-domain-mean} 定义的 $v$ 在 $\Omega$ 上的平均值. 对式~\eqref{eq:ch05-mean-zero-variational-space} 定义的 $V$, $v\in V$ 蕴含 $\overline v=0$, 故 $\|v\|_{L^2(\Omega)}^2\leq C_\Omega^2|v|_{H^1(\Omega)}^2=C_\Omega^2a(v,v)$. 两边相加 $a(v,v)$ 即得所需的强制性不等式.

\MBSourceItem{2}
\begin{Proposition}
\label{prop:ch05-pure-neumann-coercivity}
设 $\Omega$ 满足式~\eqref{eq:ch05-best-constant-poincare}, 例如 $\Omega$ 是有限个上述星形区域之并. 令 $V$ 由式~\eqref{eq:ch05-mean-zero-variational-space} 定义, $C_\Omega$ 为式~\eqref{eq:ch05-best-constant-poincare} 中的逼近常数. 则
\[
\MathBookFormulaID{formula-ch05-pure-neumann-coercivity}
\|v\|_{H^1(\Omega)}^2\leq(1+C_\Omega^2)a(v,v)
\quad\forall v\in V.
\]
\end{Proposition}

\MathBookSourcePage{147}
现在验证 Dirichlet 边界条件情形的强制性, 其中 $V$ 由式~\eqref{eq:ch05-mixed-variational-space} 定义. 先对 $H^1(\Omega)$ 建立更一般的不等式
\MBSourceItem{3}
\begin{equation}
\MathBookFormulaID{formula-ch05-boundary-poincare-inequality}
\label{eq:ch05-boundary-poincare-inequality}
\|v\|_{L^2(\Omega)}
\leq C\left(\left|\int_\Gamma v\,ds\right|+|u|_{H^1(\Omega)}\right)
\quad\forall v\in H^1(\Omega),
\end{equation}
其中正常数 $C$ 只依赖于 $\Omega$ 和 $\Gamma$. 任取 $v\in H^1(\Omega)$, 有
\[
\MathBookFormulaID{formula-ch05-boundary-poincare-proof-one}
\begin{aligned}
\|v\|_{L^2(\Omega)}
&\leq\|v-\overline v\|_{L^2(\Omega)}+\|\overline v\|_{L^2(\Omega)}\\
&\leq C_\Omega|v|_{H^1(\Omega)}
+\frac{|\Omega|^{1/2}}{|\Gamma|}\left|\int_\Gamma v\,ds\right|\\
&\leq C_\Omega|v|_{H^1(\Omega)}
+\frac{|\Omega|^{1/2}}{|\Gamma|}
\left(\left|\int_\Gamma\overline v\,ds\right|
+\left|\int_\Gamma(v-\overline v)\,ds\right|\right),
\end{aligned}
\]
并且
\[
\MathBookFormulaID{formula-ch05-boundary-poincare-proof-two}
\begin{aligned}
\left|\int_\Gamma(v-\overline v)\,ds\right|
&\leq|\Gamma|^{1/2}\|v-\overline v\|_{L^2(\partial\Omega)}\\
&\leq|\Gamma|^{1/2}C_\Omega
\bigl(\|v-\overline v\|_{L^2(\Omega)}+|v-\overline v|_{H^1(\Omega)}\bigr)\\
&\leq|\Gamma|^{1/2}C_\Omega|v|_{H^1(\Omega)}.
\end{aligned}
\]
两式合起来即得式~\eqref{eq:ch05-boundary-poincare-inequality}.

\MBSourceItem{4}
\begin{Proposition}
\label{prop:ch05-mixed-boundary-coercivity}
设 $\Omega$ 如命题~\ref{prop:ch05-pure-neumann-coercivity}, 且 $\Gamma$ 为闭集, $\operatorname{meas}(\Gamma)>0$. 令 $V$ 由式~\eqref{eq:ch05-mixed-variational-space} 定义. 则存在只依赖于 $\Omega$ 和 $\Gamma$ 的常数 $C$, 使
\[
\MathBookFormulaID{formula-ch05-mixed-boundary-coercivity}
\|v\|_{H^1(\Omega)}^2\leq Ca(v,v)\quad\forall v\in V.
\]
\end{Proposition}

这一估计在 $L^p$ 框架中也成立(参见练习~\ref{exr:ch05-13}). 特殊情形 $\Gamma=\partial\Omega$ 经常使用, 并有专名. 回忆 $\mathring W_p^1(\Omega)$ 的记号(参见式~\eqref{eq:ch01-zero-trace-w1}), 它表示满足 $\partial\Omega$ 上 Dirichlet 边界条件的 $W_p^1(\Omega)$ 函数.

\MBSourceItem{5}
\begin{Proposition}[Poincaré 不等式]
\label{prop:ch05-poincare-inequality}
设 $\Omega$ 如命题~\ref{prop:ch05-pure-neumann-coercivity}. 则存在常数 $C<\infty$, 使
\[
\MathBookFormulaID{formula-ch05-poincare-inequality}
\|v\|_{W_p^1(\Omega)}\leq C|v|_{W_p^1(\Omega)}
\quad\forall v\in\mathring W_p^1(\Omega).
\]
\end{Proposition}
